\section{Related Work}
\label{sec:related-work}

Work relevant to this paper answers four questions using four kinds of evidence. \emph{Prospective structural evaluation} asks what an attacker could achieve given a model of configurations, vulnerabilities, and reachability. \emph{Empirical replay observation} asks what a deployed defense did when specific actions were executed against it. \emph{Defense-effectiveness scoring} asks how such results combine into a comparable quantity. \emph{Repair attribution} asks which control should be fixed as a consequence. The four are not interchangeable: a reachable path, an alert, a blocking decision, and a remediation target are different objects, and a score over one cannot be read as a statement about another. We therefore organize prior work by the question it answers and by its unit of analysis.

\subsection{Attack Knowledge Bases, Adversary Emulation, and Defense Benchmark Construction}
\label{subsec:rw-emulation-benchmark}

ATT\&CK supplies the tactic and technique vocabulary used to normalize adversary behavior and derive canonical phases in our benchmark~\cite{mitreAttack}. ATT\&CK Evaluations record how individual products expose or detect procedures during emulation~\cite{mitreAttackEvaluations}, while tooling such as CALDERA~\cite{caldera} makes technique-level tests executable and harnesses such as DELTA orchestrate attacks against a target~\cite{lee2017delta}. These resources yield executable actions and per-product observations, but do not specify how such observations become comparable outcomes, nor how an aggregate points back to a control worth repairing.

Two lines of work evaluate a deployed defense at action granularity by injection or replay, and they are the closest prior art to our evaluation pipeline. SAIBERSOC injects synthetic ATT\&CK-derived attacks into a live security operations center~\cite{rosso2022saibersoc}. Its evaluated object is the operational center as a whole, monitoring infrastructure together with the analysts who staff it; its unit of analysis is an injected attack instance; and it reports detection accuracy separated into identification and investigation, together with time-to-investigation, compared between differently configured centers. Cao et al. replay real-world attack variants against several detection methods and report detection capability per variant~\cite{cao2016replay}; their evaluated object is the detector and their unit is the replayed variant. Both pipelines coincide with ours up to the moment an outcome is recorded. They diverge in what that outcome is designed to explain: SAIBERSOC explains variance attributable to center configuration and analyst process, and Cao et al. explain per-variant detection capability, each appropriate to its own evaluation objective. Our evaluation record instead carries four elements together, the benchmark action, a reference primary control-point attribution fixed before evaluation, the observed defense outcome, and the action's campaign-stage context, so that an aggregate can be decomposed along any of them.

Automated workflows spanning defense generation, enforcement, and evaluation occupy adjacent ground. CD-GEE covers these phases and validates on case-study networks~\cite{enoch2022cdgee}. Its reported evaluation unit is a difference in security posture before and after defense deployment, whereas ours is an outcome observed for a specific evaluated benchmark action. The comparison therefore operates at different granularities: system-level posture in CD-GEE and action-level defense evidence in our framework.

Mappings between controls and attacker behavior are maintained both in research and in industry. D3FEND organizes countermeasures and their relationships to adversary behavior~\cite{mitreD3FEND,kaloroumakis2021d3fend}, prospectively identifying which defensive techniques may be relevant to a behavior, and the MITRE Center for Threat-Informed Defense publishes thousands of technique-to-control mappings with a documented methodology~\cite{mitreCTIDMappings}. Both organize catalogs of defensive controls. Our framework instead uses a victim-side control-point taxonomy and a documented crosswalk between attributed control points and primary or compensating defense categories, and reports agreement statistics for a sampled annotation pilot. Vendor-specific practitioner workflows provide further context: breach-and-attack-simulation and adversarial-exposure-validation products execute attack procedures against a live environment and report control gaps to their operators. We note the practice for its shared concern with observed rather than assumed coverage, and draw no comparison with it, since its methods and results are documented per product.

Our scoring layer consumes observations from systems never designed to be scored. Provenance-based intrusion detection has been systematized as a field~\cite{inam2023sokprovenance}; representatively, HOLMES correlates suspicious flows into campaign-level scenarios~\cite{milajerdi2019holmes} and KAIROS reconstructs compact attack summaries from whole-system provenance~\cite{cheng2024kairos}. On the enforcement side, SANE renders some connectivity architecturally impossible~\cite{casado2006sane}, FIREMAN yields static policy findings~\cite{yuan2006fireman}, and programmable-data-plane defenses interrupt traffic during execution~\cite{zhang2020poseidon,yoo2024smartcookie}. These are representatives of evidence types rather than an inventory: what matters for scoring is the distinction between architectural impossibility, runtime interruption, observation without interruption, and absence of evidence, which map to \texttt{PREVENTED}, \texttt{BLOCKED}, \texttt{DETECTED}, and \texttt{MISSED}.

\subsection{ATT\&CK-Based Coverage and Evaluation Reinterpretation}
\label{subsec:rw-attck-coverage}

A coverage count over ATT\&CK techniques is a tempting capability metric, and it has already been tested. Virkud et al. examine how endpoint detection products and public rulesets use ATT\&CK and show empirically that technique coverage is a poor indicator of defensive capability~\cite{virkud2024endpoint}. Their result is diagnostic, and we build on it rather than re-deriving it. We therefore do not use technique coverage as the score. Evaluation is operationalized one level down, at an attributed action: attributed benchmark actions carry a reference primary control-point attribution, and it is the observed outcome on such an action, rather than the presence of a technique in a coverage list, that enters the score.

Step-level results have also been aggregated across an attack chain. Shen et al. re-link the steps of a published ATT\&CK Evaluations round into causal attack graphs, motivated by the observation that published reports present steps in isolation~\cite{shen2024decoding}. The inputs differ in direction: their evidence is the released output of one existing evaluation, scoped to endpoint products, while ours is a benchmark and attribution layer built from playbooks, which allows a previously unevaluated cross-category defense stack to be scored and its score loss traced to a control.

Both depend on the provenance of the mappings that produce them. Threat-intelligence extraction~\cite{liao2016ioc} and temporal-relation learning over attack steps~\cite{shen2019attack2vec} show such mappings to be derived artifacts, and operational feeds differ substantially in coverage and timeliness~\cite{li2019tealeaves}. We therefore record a confidence label with each attributed action, and the coding guide used during annotation distinguishes how an attribution was reached, rather than treating all mappings as equivalent evidence.

\subsection{Defense Capability and Effectiveness Scoring}
\label{subsec:rw-scoring}

Comparable defense evaluation has been unified before, on a different denominator. Iannacone and Bridges survey the area and express intrusion-detection metrics, cyber-exercise scoring, and security economics through a common cost term~\cite{iannacone2020quantifiable}. Ours is a stage-weighted interception value paired with a repairable control target, so the output is a score that localizes rather than a cost that totals. The cost-oriented lineage is longer still~\cite{butler2002saem,gordon2002economics}, and it illustrates a failure mode we avoid: collapsing incomparable consequences into one index produces a number that is hard to act on, which is why we also report by phase, control-point domain, and product category.

DEFIA is particularly close in intent, scoring the utility of attack behaviors prevented and recording the point at which a defense acts~\cite{liu2023defia}. Its unit of analysis is an attack behavior and its semantics are effectiveness-oriented, resting on whether a behavior was prevented and on where the defense took effect. Three properties of our unit separate the two. Observed outcomes occupy a four-level normalized scale, \texttt{PREVENTED}, \texttt{BLOCKED}, \texttt{DETECTED}, and \texttt{MISSED}, rather than a prevented/not-prevented distinction. Attributed benchmark actions carry a reference primary control-point attribution assigned during benchmark construction, independently of the system under test. And stage membership carries explicit alternative, conjunctive, and clustered composition, which our benchmark must retain when aggregating observed action outcomes from action to stage to playbook.

Prior work on moving-target-defense evaluation already separates complementary views of defense effectiveness. Zaffarano et al. define a quantitative effectiveness framework in which mission-side and attack-side outcomes are measured separately rather than folded into a single figure~\cite{zaffarano2015mtd}, and Taylor et al. apply that framework under automated experimentation, reporting the mission and attack metrics as a pair~\cite{taylor2016mtd}. This separation is the closest structural precedent for our coverage-versus-hit-rate reporting. The evaluated object differs: theirs is a single defense class, moving-target mechanisms, whereas ours is a heterogeneous multi-product stack, so the pairing becomes coverage of the action set structurally applicable to a product category, measured against the hit rate observed within that coverage. Neither work attaches a victim-side control reference to the action or composes outcomes under within-stage semantics. Detection outcomes have likewise been weighted by time: Garcia et al. apply a correcting function per window so that earlier detection scores higher~\cite{garcia2014botnet}. The weighted quantity differs, elapsed detection time in their case and semantic attack phase combined with stage dependency in ours. Time-sensitive and multi-dimensional weighted scoring also appears in live exercises~\cite{doupe2011hitem,werther2011ctf}, where the measured object is a participating team.

Dong et al. assess provenance-based endpoint detection and response tools against operational decision factors such as client overhead and triage cost~\cite{dong2023arewethereyet}. Those dimensions measure the cost of operating a tool; ours characterize observed defense outcomes in campaign context while retaining benchmark-side control-point attribution.

\subsection{Multi-Stage Evaluation, Attack-Path Reasoning, and Repair Prioritization}
\label{subsec:rw-multistage-repair}

Structural methods evaluate multi-stage risk without observing a defense in operation. CVSS scores the severity of an individual vulnerability~\cite{mell2007common,first2019cvss31}, and attack-graph methods reason over configurations, privileges, and reachability through symbolic model checking~\cite{sheyner2002attackgraphs}, logical derivation~\cite{ou2005mulval}, or generation at larger scale~\cite{ou2006scalable}. Derived metrics extend the same view: $k$-zero-day safety counts the unknown vulnerabilities an attacker would require~\cite{wang2014kzero}, network diversity estimates resilience from software heterogeneity~\cite{zhang2016networkdiversity}, and mean time-to-compromise estimates attacker progress across protection zones~\cite{leversage2008mttc}. Each is computed from a model of what is reachable rather than from what a deployed defense did. Attack-graph-based security measurement composes such models into overall metrics for a networked system~\cite{wang2007measuring}, and system-level assessment aggregates risk across an organization's assets~\cite{erola2022system}. These lines organize progression, feasibility, and system-level risk, and they are well suited to those questions. The dimension they organize is nonetheless not the one our attribution supplies: where an action sits in a campaign, and how far an attacker could progress, is a separate question from which class of victim-side control the benchmark associates with that action. We therefore keep campaign phase and control-point domain as independent fields rather than deriving either from the other. Our crosswalk uses structural information in the same prospective spirit, to decide which product categories are expected to cover an action, and reports that applicability separately from empirical performance.

Attack steps have been assigned differentiated defensive value in this setting. GPLADD models multi-step attacker-defender interaction with explicit per-step timing and value~\cite{outkin2019gpladd}, and later work applies the approach to defender policy evaluation and resource allocation using ATT\&CK Evaluations data~\cite{outkin2022defender}. The mode of inference separates the two: that work optimizes allocation probabilistically under uncertainty and outputs a policy, whereas ours computes a deterministic score from observed outcomes and outputs a figure reducible to the attributed actions that account for each score loss.

Repair targets have long been selected from attack models, and this is the prior-work line closest to our attribution contribution. Noel et al. reason over an exploit-dependency graph to identify the initial conditions whose removal makes critical resources unreachable, returning a hardening set rather than a per-event judgement~\cite{noel2003hardening}; Wang, Noel and Jajodia develop the minimum-cost formulation of the same problem, selecting a remediation configuration under cost constraints~\cite{wang2006hardening}. Both operate prospectively over a system and attack model, and both output a set chosen to be sufficient for a reachability goal. The overlap is real: each produces a remediation target from an attack representation. The difference lies in when the target is fixed and what it is attached to. The framework records benchmark-side primary control-point references during benchmark construction, before any system is evaluated, and subsequently aligns observed outcomes to the corresponding benchmark records, aggregating them while action, stage, and control-point traceability is preserved. Unlike minimum-cost hardening, it does not solve for an optimal or sufficient remediation set. Neither output subsumes the other.

The word \emph{attribution} requires disambiguation, because it already denotes a different object in this literature. ORTHRUS defines quality of attribution as the analyst effort needed to recover an attack's causal origin from detection output~\cite{jiang2025orthrus}, and R-CAID embeds root-cause analysis within provenance-based detection~\cite{goyal2024rcaid}. Both reconstruct, explain, or localize attacker-side causal origins in provenance and event space, and both are evaluated on how well that reconstruction succeeds. The attribution in this paper is a different object on a different substrate: a benchmark-side victim-side control-point reference attached to a benchmark action, used for defense scoring and repair-oriented decomposition rather than for explaining an observed intrusion. Neither output subsumes the other, and the shared word should not be read as a shared task.

\begin{table*}[!t]
\centering
\small
\caption{Positioning by evidence, unit of analysis, output, and remediation traceability. Rows group research lines by the question each answers; the grouping does not indicate competitor class.}
\label{tab:related-work-position}
\setlength{\tabcolsep}{3pt}
\begin{tabular}{@{}p{0.155\textwidth}p{0.185\textwidth}p{0.145\textwidth}p{0.190\textwidth}p{0.165\textwidth}@{}}
\hline
\textbf{Research line} & \textbf{Primary evidence} & \textbf{Unit of analysis} & \textbf{Typical output} & \textbf{Remediation traceability} \\
\hline
Knowledge bases, emulation, and replay & Technique descriptions and emulation or replay telemetry & Technique, emulated step, or injected attack & Vocabulary, executable test, or per-attack observation & Partial; process- or capability-oriented \\
ATT\&CK coverage and reinterpretation & Published evaluation outputs and detection rulesets & Technique or evaluation step & Coverage assessment or re-linked attack graph & Indirect; scoped to one product class \\
Effectiveness scoring & Costs, mission metrics, or prevented-behavior utility & Cost-accruing event or attack behavior & Comparable effectiveness or cost index & Aggregate; not tied to one control \\
Multi-stage reasoning and repair prioritization & Configuration, topology, privileges, and modelled paths & Vulnerability, condition, or attack path & Severity, reachability, or minimal hardening set & Structural and prospective \\
Validity and reproducibility & Published benchmarks and evaluation designs & Benchmark, metric, or replication attempt & Methodological constraint or checklist & Not applicable \\
Industry mappings and practitioner validation & Control catalogs and executed attack procedures & Technique--control pair or executed procedure & Mapping set or vendor gap report & Control-oriented; documented per product \\
This work & Playbook action and aligned defense observation & Action--outcome--control record & Stage-aware, multi-dimensional defense score & Explicit reference primary control-point attribution \\
\hline
\end{tabular}
\end{table*}

\subsection{Evaluation Validity, Reproducibility, and Operational Evidence}
\label{subsec:rw-validity}

A unified layer over heterogeneous systems is itself a known construction. Bilot et al. build one for provenance-based intrusion detection, ingesting multiple implementations and reporting comparable metrics~\cite{bilot2025simpler}. The evaluated object differs: detection quality of research prototypes on public datasets in their case, a defense stack evaluated against benchmark playbook actions in our framework, where attributed benchmark records carry a reference primary control-point attribution.

Several results constrain how our scores may be interpreted. Sommer and Paxson argue that designing the evaluation is harder than building the detector~\cite{sommer2010closedworld}, and the base-rate fallacy shows that a low false-positive rate can coexist with a high false-alarm rate under class imbalance~\cite{axelsson2000baserate}, which is why detection denominators must be kept apart. We apply that separation to product-category scoring, reporting coverage over the structurally applicable action set apart from hit rate within it. Verendel finds that quantified-security methods have generally lacked repeated testing~\cite{verendel2009quantified}, and Herley and van Oorschot argue that security claims require explicit assumptions, evidence, and limits of inference~\cite{herley2017science}. We follow these constraints by stating the evidence basis of each reported quantity.

Benchmark validity should be assessed rather than assumed. Van der Kouwe et al. catalog recurring flaws in systems benchmarking and motivate treating benchmark-validity limitations as an explicit subject rather than an implicit one~\cite{vanderkouwe2019benchmarkingflaws}, while Olszewski et al. show that reproduction claims depend on environment and implementation assumptions being made explicit~\cite{olszewski2025replicability}. Both concerns apply to a benchmark of this kind and shape how its reported quantities should be qualified.

One scope limitation follows. The attribution schema covers endpoint, identity, data, network, and operational controls, so a network-based testbed can validate only a replayable slice of the schema rather than its full scope. Actions requiring unsupported endpoint state should be treated as out of scope or as explicit preconditions rather than being assigned a defense outcome.

\subsection{Synthesis and Positioning}
\label{subsec:rw-position}

Table~\ref{tab:related-work-position} summarizes these lines by evidence, unit of analysis, output, and remediation traceability. Existing work separately contributes attacker and action vocabularies together with replay evidence, defense effectiveness and capability metrics, multi-stage and attack-path reasoning, and repair-prioritization or causal-attribution outputs. Each of these components is established, and we claim none of them in isolation. What this framework combines is a benchmark-side reference primary control-point attribution, an observed action outcome on a four-level normalized scale, and explicit within-stage composition semantics, held together so that a comparable defense score remains traceable to the attributed actions and control points that account for it. Simply composing the existing outputs would not achieve this, because an alert, a blocking decision, a posture difference, and a prospective path-risk value do not identify the same evaluated event and cannot be aggregated without an explicit alignment rule. That alignment rule, and the traceability it preserves, is what the action--outcome--control record provides.
